\documentclass[a4paper,11pt]{article}
\usepackage[utf8]{inputenc}
\usepackage[T1]{fontenc}
\usepackage[ngerman]{babel}
\usepackage{geometry}
\geometry{left=2.5cm, right=2.5cm, top=2.5cm, bottom=2.5cm}
\usepackage{longtable}
\usepackage{booktabs}
\usepackage{array}
\usepackage{xcolor}
\usepackage{colortbl}
\usepackage{hyperref}
\usepackage{fancyhdr}
\usepackage{listings}
\usepackage{tabularx}
\usepackage{enumitem}

\definecolor{headerblue}{HTML}{1565C0}
\definecolor{tableheader}{HTML}{E3F2FD}
\definecolor{codebg}{HTML}{F5F5F5}
\definecolor{codegreen}{HTML}{2E7D32}
\definecolor{codegray}{HTML}{757575}

\lstset{
  backgroundcolor=\color{codebg},
  basicstyle=\ttfamily\small,
  keywordstyle=\color{headerblue}\bfseries,
  commentstyle=\color{codegreen},
  stringstyle=\color{codegray},
  breaklines=true,
  frame=single,
  rulecolor=\color{codegray!30},
  numbers=none,
  showstringspaces=false,
  tabsize=4,
  language=C++
}

\hypersetup{colorlinks=true,linkcolor=headerblue,urlcolor=headerblue}

\pagestyle{fancy}
\fancyhf{}
\fancyhead[L]{\textbf{INSPECTAIR}}
\fancyhead[R]{Coding Rules v2.0}
\fancyfoot[C]{\thepage}
\renewcommand{\headrulewidth}{0.4pt}

\setlength{\parindent}{0pt}
\setlength{\parskip}{6pt}

\begin{document}

\begin{center}
{\LARGE\textbf{INSPECTAIR}}\\[0.3cm]
{\Large Coding Rules}\\[0.3cm]
{\large Programmierrichtlinien und Konventionen}\\[0.5cm]
\begin{tabular}{ll}
\textbf{Projekt:} & InspectAir -- Luftqualitätsmessgerät \\
\textbf{Version:} & 2.0 \\
\textbf{Datum:} & Februar 2026 \\
\textbf{Autoren:} & Maximilian Liebl, Cedric Stroh, Jannis Messer \\
\textbf{Modul:} & Mikrocomputertechnik 1, DHBW Ravensburg \\
\textbf{Dozent:} & Thomas Stingl (Airbus) \\
\end{tabular}
\end{center}

\vspace{0.5cm}
\noindent\textit{Dieses Dokument wurde unter Verwendung von KI-Tools (Claude, Anthropic) zur Überprüfung erstellt.}
\vspace{0.3cm}

\tableofcontents
\newpage

% ============================================================
\section{Programmiersprache und Umgebung}

\begin{tabularx}{\textwidth}{lX}
\toprule
\rowcolor{tableheader}
\textbf{Parameter} & \textbf{Wert} \\
\midrule
Sprache & C++ (Arduino Framework) \\
IDE & VSCode mit PlatformIO \\
Board & esp32-s3-devkitc-1 \\
Framework & Arduino (Espressif32) \\
Flash Mode & QIO 80\,MHz, OPI PSRAM \\
Upload Speed & 921600 Baud \\
Monitor Speed & 115200 Baud \\
\bottomrule
\end{tabularx}

% ============================================================
\section{Formatierung}

\begin{tabularx}{\textwidth}{lX}
\toprule
\rowcolor{tableheader}
\textbf{Regel} & \textbf{Beschreibung} \\
\midrule
Einrückung & 4 Leerzeichen (keine Tabs) \\
Zeilenlänge & Maximal 100 Zeichen \\
Encoding & UTF-8 \\
Zeilenende & LF (Unix-Style) \\
Klammern & Öffnende geschweifte Klammer auf gleicher Zeile (K\&R-Stil) \\
Leerzeilen & Eine Leerzeile zwischen logischen Blöcken und Funktionen \\
\bottomrule
\end{tabularx}

% ============================================================
\section{Namenskonventionen}

\begin{tabularx}{\textwidth}{llX}
\toprule
\rowcolor{tableheader}
\textbf{Element} & \textbf{Konvention} & \textbf{Beispiel} \\
\midrule
Variablen & \texttt{snake\_case} & \texttt{sensor\_data}, \texttt{last\_read\_time} \\
Funktionen & \texttt{snake\_case} & \texttt{read\_sensors()}, \texttt{ui\_init()} \\
Konstanten / Defines & \texttt{UPPER\_SNAKE\_CASE} & \texttt{PRESENCE\_TIMEOUT\_DIM\_MS} \\
Structs / Enums & \texttt{PascalCase} & \texttt{SensorReadings}, \texttt{DisplayState} \\
Enum-Werte & \texttt{UPPER\_SNAKE\_CASE} & \texttt{UI\_SCREEN\_OVERVIEW}, \texttt{STATE\_ACTIVE} \\
Klassen & \texttt{PascalCase} & \texttt{PowerManager}, \texttt{WifiClock} \\
Private Member & Präfix \texttt{\_} & \texttt{\_state}, \texttt{\_fadeCurrent} \\
Typedefs & Suffix \texttt{\_Data} & \texttt{AHT20\_Data}, \texttt{PMS5003\_Data} \\
Header Guards & \texttt{UPPER\_SNAKE\_CASE} mit \texttt{\_H} & \texttt{POWER\_MANAGER\_H} \\
Dateinamen & \texttt{snake\_case.cpp/.h} & \texttt{power\_manager.cpp}, \texttt{sensor\_filter.h} \\
Pin-Defines & Präfix nach Funktion & \texttt{TFT\_MOSI}, \texttt{I2C\_SDA}, \texttt{CO2\_RX} \\
\bottomrule
\end{tabularx}

% ============================================================
\section{Dokumentation}

Alle öffentlichen Funktionen, Klassen und Strukturen werden mit Doxygen-Kommentaren versehen:

\begin{lstlisting}
/**
 * @brief Aktualisiert alle Sensorwerte auf dem Display
 * @param readings Aktuelle Sensor-Messwerte
 */
void ui_updateSensorValues(const SensorReadings& readings);
\end{lstlisting}

Inline-Kommentare erklären nicht-offensichtliche Logik:
\begin{lstlisting}
// 220uF Stuetzkondensator puffert Stromspitzen des MH-Z19C
delay(100);  // Warten bis Sensor-Versorgung stabil
\end{lstlisting}

% ============================================================
\section{Fehlerbehandlung}

Da Exceptions auf ESP32 nicht verfügbar sind, wird folgendes Schema verwendet:

\begin{itemize}[nosep]
\item Jeder Sensor-Treiber gibt ein \texttt{ok}-Flag in seiner Datenstruktur zurück (z.\,B. \texttt{MHZ19C\_Data.valid}).
\item Fehlerzustände werden über \texttt{Serial.println()} auf der Debug-Konsole protokolliert.
\item Die \texttt{main.cpp} prüft bei jedem Lesezyklus den Rückgabewert und zeigt fehlerhafte Sensoren auf dem Display an.
\item Ein Watchdog-Timer (WDT) überwacht die Hauptschleife und löst bei Blockierung einen Reset aus.
\end{itemize}

Beispiel aus dem Sensor-Treiber:
\begin{lstlisting}
MHZ19C_Data mhz19c_read() {
    MHZ19C_Data data;
    data.co2_ppm = myMHZ19.getCO2();
    data.valid = (data.co2_ppm > 0 && data.co2_ppm < 5000);
    return data;
}
\end{lstlisting}

% ============================================================
\section{Verbotene Praktiken}

\begin{tabularx}{\textwidth}{lX}
\toprule
\rowcolor{tableheader}
\textbf{Verboten} & \textbf{Begründung} \\
\midrule
\texttt{delay()} in \texttt{loop()} & Blockiert LVGL-Rendering und Sensor-Abfragen. Stattdessen \texttt{millis()}-basiertes Timing verwenden. \\
\texttt{String} Klasse & Heap-Fragmentierung auf ESP32. Stattdessen \texttt{char[]}, \texttt{snprintf()} oder \texttt{lv\_label\_set\_text\_fmt()} verwenden. \\
Globale Variablen & Nur \texttt{static}-Variablen in \texttt{main.cpp} und \texttt{extern}-Singleton-Instanzen (z.\,B. \texttt{powerManager}) erlaubt. \\
Magic Numbers & Alle Konstanten als \texttt{\#define} oder \texttt{static const} mit sprechendem Namen in \texttt{pins.h}, \texttt{colors.h} oder \texttt{app\_config.h}. \\
Dynamische Allokation in \texttt{loop()} & Kein \texttt{new}/\texttt{malloc} in der Hauptschleife. LVGL-Objekte werden einmalig in \texttt{ui\_init()} erstellt. \\
Blocking I/O & Sensor-Reads mit Timeout absichern; niemals endlos auf Daten warten. \\
\bottomrule
\end{tabularx}

% ============================================================
\section{Versionskontrolle}

\begin{tabularx}{\textwidth}{lX}
\toprule
\rowcolor{tableheader}
\textbf{Regel} & \textbf{Beschreibung} \\
\midrule
Repository & GitHub (privat), Git-basiert \\
Branching & \texttt{main} (stabil) + Feature-Branches (\texttt{feature/sensor-xyz}) \\
Commit-Format & Conventional Commits: \texttt{feat:}, \texttt{fix:}, \texttt{docs:}, \texttt{refactor:} \\
Merge-Strategie & Feature-Branch $\to$ Pull Request $\to$ Review $\to$ Merge in \texttt{main} \\
.gitignore & Build-Artefakte (\texttt{.pio/}), IDE-Konfiguration (\texttt{.vscode/}), Credentials \\
\bottomrule
\end{tabularx}

% ============================================================
\section{Build-Konfiguration}

Zentrale Build-Flags in \texttt{platformio.ini}:

\begin{lstlisting}
build_flags =
  -DARDUINO_USB_CDC_ON_BOOT=1
  -DBOARD_HAS_PSRAM
  -DLV_CONF_INCLUDE_SIMPLE
  -DLV_COLOR_DEPTH=16
  -DLV_COLOR_16_SWAP=1
  -DLV_MEM_SIZE="(96U * 1024U)"
  -DLV_FONT_MONTSERRAT_16=1
\end{lstlisting}

Pin-Definitionen sind zentral in \texttt{pins.h} gesammelt und dürfen nicht in einzelnen Modulen hart codiert werden.

% ============================================================
\section{Änderungshistorie}

\begin{tabularx}{\textwidth}{l l l X}
\toprule
\rowcolor{tableheader}
\textbf{Version} & \textbf{Datum} & \textbf{Autor} & \textbf{Änderung} \\
\midrule
1.0 & Januar 2026 & Team & Erstversion \\
2.0 & Februar 2026 & Team & Ergänzung: Watchdog, LVGL-spezifische Regeln, Build-Flags, Private-Member-Konvention, Dateinamen-Konvention. \\
\bottomrule
\end{tabularx}

\end{document}
